\documentclass[12pt,a4paper]{article}

\usepackage[utf8]{inputenc}
\usepackage[T1]{fontenc}
\usepackage[brazil]{babel}
\usepackage{geometry}
\usepackage{setspace}
\usepackage{booktabs}
\usepackage{array}
\usepackage{hyperref}
\usepackage{listings}
\usepackage{xcolor}

\geometry{left=3cm,right=2cm,top=3cm,bottom=2cm}
\onehalfspacing

\lstset{
  basicstyle=\ttfamily\small,
  breaklines=true,
  frame=single,
  backgroundcolor=\color{gray!8},
  columns=fullflexible
}

\title{
  \textbf{Implementação de uma Rede P2P para Simulação de DTN em Cenários de Alta Latência e Desconexão}
}
\author{
  Universidade Federal do Amazonas\\
  Disciplina de Sistemas Distribuídos\\[0.5cm]
  Márcio Éric Lamêgo Valente\\
  José Santo de Moura Neto
}
\date{\today}

\begin{document}

\maketitle
\thispagestyle{empty}

\newpage
\tableofcontents
\newpage

\section{Introdução}

Este relatório apresenta a implementação de uma rede P2P simulada para cenários de DTN
(\textit{Delay-Tolerant Network}), conforme a atividade prática da disciplina de Sistemas
Distribuídos. O objetivo principal foi demonstrar, por meio de simulação, como mensagens
podem ser armazenadas e encaminhadas em uma rede sem conectividade fim-a-fim garantida.

Em redes DTN, os nós podem ficar desconectados durante longos períodos. Por isso, cada nó
atua como uma entidade autônoma capaz de armazenar mensagens localmente e encaminhá-las
quando encontra outro nó. Esse modelo é conhecido como \textit{store-and-forward}.

\section{Objetivos}

Os objetivos da implementação foram:

\begin{itemize}
  \item simular uma arquitetura P2P sem servidor central;
  \item representar cada nó com identificador, buffer de mensagens, relógio lógico e lista de mensagens vistas;
  \item implementar roteamento epidêmico para encaminhamento de mensagens;
  \item simular contatos intermitentes entre pares de nós;
  \item avaliar taxa de entrega, latência média e overhead;
  \item executar testes em cenários denso, esparso e em cadeia.
\end{itemize}

\section{Arquitetura da Solução}

A solução foi implementada em Python e utiliza uma simulação baseada em eventos discretos.
Cada nó P2P é representado por uma entidade independente dentro do simulador. O simulador
não funciona como servidor central de mensagens; ele apenas controla o avanço do tempo e
agenda os eventos de contato definidos nos arquivos de entrada.

Os principais módulos do sistema são:

\begin{table}[h]
\centering
\begin{tabular}{>{\raggedright\arraybackslash}p{4cm} >{\raggedright\arraybackslash}p{9cm}}
\toprule
\textbf{Módulo} & \textbf{Responsabilidade} \\
\midrule
\texttt{SimulationEngine} & Carrega entradas, agenda eventos e coordena a simulação. \\
\texttt{ContactParser} & Lê o arquivo de contatos no formato \texttt{tempo\_inicio, tempo\_fim, noA, noB}. \\
\texttt{MessageParser} & Lê mensagens iniciais e gera identificadores de mensagens. \\
\texttt{Node} & Representa um nó P2P com buffer, mensagens vistas, relógio lógico e probabilidades de entrega. \\
\texttt{EpidemicRouter} & Implementa a troca de resumos e transferência de mensagens ausentes. \\
\texttt{MetricsCollector} & Registra entregas, transmissões, latência e overhead. \\
\texttt{CLI} & Disponibiliza comandos para execução da simulação. \\
\bottomrule
\end{tabular}
\end{table}

\section{Modelo de Contatos}

Os contatos são definidos em arquivos \texttt{contacts.txt}. Cada linha indica o intervalo
em que dois nós podem se comunicar:

\begin{lstlisting}
tempo_inicio, tempo_fim, noA, noB
\end{lstlisting}

Exemplo:

\begin{lstlisting}
0.0, 5.0, 1, 2
5.1, 8.0, 2, 3
8.2, 12.0, 3, 1
\end{lstlisting}

Durante a simulação, os eventos são processados em ordem crescente de tempo. Quando ocorre
um contato, os dois nós envolvidos executam a troca P2P.

\section{Roteamento Epidêmico}

O roteamento epidêmico foi implementado com base na troca de resumos de mensagens. Quando
dois nós entram em contato, cada um informa ao outro quais mensagens já conhece. Em seguida,
cada nó envia as mensagens que possui e que o outro ainda não viu.

O processo executado em cada contato é:

\begin{enumerate}
  \item atualizar o relógio lógico dos nós para o tempo do contato;
  \item obter o conjunto de mensagens vistas por cada nó;
  \item identificar as mensagens que o outro nó ainda não possui;
  \item filtrar mensagens com TTL maior que zero;
  \item copiar a mensagem para o receptor, decrementando o TTL;
  \item registrar a transmissão;
  \item registrar entrega caso o receptor seja o destino final.
\end{enumerate}

O uso da lista de mensagens vistas evita loops, pois um nó não aceita novamente uma mensagem
que já recebeu anteriormente. O TTL limita a quantidade máxima de saltos de cada mensagem,
evitando propagação indefinida.

\section{Interface de Execução}

A execução principal é feita por linha de comando. Um exemplo de execução do cenário em
cadeia com roteamento epidêmico é:

\begin{lstlisting}
python3 dtp2p.py simulate \
  --contacts docs/cenarios/cadeia/contacts.txt \
  --messages docs/cenarios/cadeia/messages.txt \
  --router epidemic \
  --log-level info
\end{lstlisting}

Também foi implementado suporte ao roteamento PROPHET como extensão:

\begin{lstlisting}
python3 dtp2p.py simulate \
  --contacts docs/cenarios/denso/contacts.txt \
  --messages docs/cenarios/denso/messages.txt \
  --router prophet
\end{lstlisting}

\section{Cenários de Teste}

Foram preparados três cenários principais:

\begin{itemize}
  \item \textbf{Cenário denso}: possui muitos contatos e múltiplos caminhos alternativos;
  \item \textbf{Cenário esparso}: possui poucos contatos e longos intervalos sem conectividade;
  \item \textbf{Cenário em cadeia}: os nós se conectam em sequência, exigindo encaminhamento por intermediários.
\end{itemize}

Esses cenários permitem observar o comportamento da rede em diferentes níveis de conectividade.
No cenário denso, espera-se alta taxa de entrega e menor latência. No cenário esparso,
espera-se maior latência e possibilidade de mensagens não entregues. No cenário em cadeia,
verifica-se diretamente o funcionamento do modelo \textit{store-and-forward}.

\section{Métricas}

As métricas avaliadas foram:

\begin{itemize}
  \item \textbf{Taxa de entrega}: razão entre mensagens entregues e mensagens criadas;
  \item \textbf{Latência média}: tempo médio entre criação e entrega das mensagens entregues;
  \item \textbf{Overhead}: razão entre o total de transmissões e o total de entregas únicas.
\end{itemize}

\section{Resultados Obtidos}

A Tabela~\ref{tab:resultados} apresenta os resultados obtidos com o roteamento epidêmico.

\begin{table}[h]
\centering
\caption{Resultados dos cenários com roteamento epidêmico}
\label{tab:resultados}
\begin{tabular}{lrrrrr}
\toprule
\textbf{Cenário} & \textbf{Mensagens} & \textbf{Entregas} & \textbf{Taxa} & \textbf{Latência média} & \textbf{Overhead} \\
\midrule
Denso & 4 & 4 & 100,00\% & 1,70s & 3,00 \\
Esparso & 4 & 3 & 75,00\% & 22,33s & 2,33 \\
Cadeia & 3 & 2 & 66,67\% & 12,75s & 4,00 \\
\bottomrule
\end{tabular}
\end{table}

Os resultados confirmam o comportamento esperado. O cenário denso obteve entrega completa,
pois havia maior disponibilidade de contatos e caminhos alternativos. O cenário esparso
apresentou menor taxa de entrega e maior latência, evidenciando o impacto da baixa
conectividade. O cenário em cadeia demonstrou o encaminhamento por nós intermediários,
caracterizando o funcionamento do modelo DTN.

\section{Testes}

A implementação possui testes automatizados usando \texttt{unittest}. A suíte cobre parsing
de contatos, parsing de mensagens, simulação por eventos, TTL, cenários principais e a
extensão PROPHET.

O comando utilizado para validação foi:

\begin{lstlisting}
python3 -m unittest discover -s tests -v
\end{lstlisting}

Todos os testes executados passaram com sucesso.

\section{Extensão PROPHET}

Além do roteamento epidêmico, foi implementada uma extensão com o roteamento PROPHET
(\textit{Probabilistic Routing Protocol using History of Encounters and Transitivity}).
Esse algoritmo utiliza probabilidades de entrega baseadas no histórico de encontros entre
nós, buscando reduzir transmissões desnecessárias quando comparado ao roteamento epidêmico.

No cenário denso, por exemplo, a execução com PROPHET entregou todas as mensagens, com
latência média de 1,95s e overhead de 2,50, indicando redução de transmissões em relação
ao roteamento epidêmico no mesmo cenário.

\section{Conclusão}

A atividade foi implementada com sucesso. O sistema simula uma rede P2P em ambiente DTN,
sem servidor central de entrega, utilizando contatos intermitentes, buffers por nó,
controle de mensagens vistas, TTL e roteamento epidêmico.

Os cenários testados demonstram a influência da conectividade na taxa de entrega, na
latência e no overhead. A implementação atende aos requisitos mínimos da atividade e ainda
inclui uma extensão com PROPHET para comparação com o roteamento epidêmico.

\end{document}
